\begin{afterword}
\summaryhead

The post reports a study in which three groups were asked how much they would pay to save 2,000, 20,000 or 200,000 migrating birds from drowning in oil ponds. They answered \$80, \$78 and \$88. The post calls this scope insensitivity: the number of birds saved barely changes what people will pay. Similar results are cited for polluted lakes in Ontario and wilderness areas in the western United States.

Two explanations are offered. In valuation by prototype, people picture one oiled bird, and the feeling that image produces sets the amount. In purchase of moral satisfaction, people pay for a warm glow whose price has nothing to do with the number of birds.

The post then reports the same pattern for human lives. A 600-fold rise in the risk from chlorinated water raised payments only fourfold, and a paper on ``psychophysical numbing'' found that saving a fixed number of lives was valued more when fewer lives were at risk. The moral is that an effective altruist must reason with the numbers, not with the image.

\responsehead

The post reports its central study correctly. The answers to the bird question barely move across a hundredfold change in the number of birds. What the post does not establish is what the answers should have been, or that the pattern carries over from survey answers to real giving.

It gives no standard for a correct response. It counts as the same error a flat response for the birds and a fourfold rise for a 600-fold rise in risk, without saying how large a rise would have been right. The question is not idle. Respondents in the bird study were told that all three numbers were 2 percent or less of the 8.5 million waterfowl in the flyway, and the post leaves this out. Someone who cares about the population has less reason to pay a hundred times more for the largest number. The moral, to think with ``the part of your brain that processes those unexciting inky zeroes,'' does not say what that part should conclude.

Several pieces of evidence are presented as firmer than they are. The prototype explanation is stated as fact before it is called a hypothesis. The quotation from Kahneman, Ritov and Schkade gains the word ``single'' and loses ``probably'' and ``perhaps.'' ``The usual finding'' of linear increases has no source, and the post's own three-level study does not show a steady increase. The paper on psychophysical numbing is said to have found that perception of deaths follows Weber's law; it mentions that law only as background. Its refugee-camp result, the post's example, reversed direction in three replications published in 2021, long after the post.

Finally, every study cited records what people say they would pay, or how they judge proposed programs. Kahneman, Ritov and Schkade, whom the post cites, argue that such dollar answers express attitudes rather than economic preferences. The bird study's own authors concluded that survey estimates fail to reflect real differences in resources, and said they had not determined why. The post draws from these surveys a moral about how to give.

In short: a correctly reported survey result, a changed quotation, no statement of what the right answer would be, and a moral about giving drawn from answers to hypothetical questions.
\end{afterword}
